\chapter{Quiet and informative markets}\label{ch:coexistence}

\begin{chapterplan}
Against one incumbent, can the price stay silent in one equilibrium and recruit a high-cost challenger in another?
This chapter proves Proposition~3 with three lemmas, and Corollary~2 locates coexistence on the example's path.
\end{chapterplan}

% =====================================================================
\section{Two ways for one market to work}\label{sec:coexistence-words}

If nobody expects the price to recruit, entry is \(\rho\), an informed order earns little, and the silent price
confirms itself: a \emph{quiet market}. If everybody expects a purchase to recruit, class \(H\) buys, and a purchase
does recruit: an \emph{informative market}, an equilibrium with informed trade.

\begin{keyidea}
Recruitment raises the return of a purchase, and informed purchases make recruitment happen.
\end{keyidea}

\paragraph{Notation.} No equilibrium has a wrong-signed order (Lemma~A.3(a)). Write \(s_+=\lambda\sigma_H(1)\),
\(s_-=\lambda\sigma_L(-1)\), and
\[
1-\mu_+(s)=\frac{(1-\pi)b}{b+\pi s},\quad \mu_-(s)=\frac{\pi b}{b+(1-\pi)s},\quad
J=\DT(1-\tau),\quad \DT^{\rm F}=\frac{k}{1-M},\quad \DT^{\rm N}=\frac{k}{\rho\max\{\pi,1-\pi\}}.
\]
So \(\mu_+\) rises from \(\pi\) to \(M\), \(\mu_-\) falls from \(\pi\) to \(m\), and \(\mu_+(s^*)=\tau\) at
\(s^*=b(\tau-\pi)/[\pi(1-\tau)]\). The returns are \(K_H=e(1)\DT[1-\mu(1)]\) and \(K_L=e(-1)\DT\mu(-1)\).

% =====================================================================
\section{Proposition 3}\label{sec:coexistence-statement}

\begin{paperresult}{Proposition 3 (Quiet and informative markets)}
Let the classes satisfy the class ranking~(1), let \(0<\pi<1\), \(0<\rho<1\), \(k>0\) and a slot \((\lambda,\beta)\)
be given, and let an indifferent investor mix over its best orders and an indifferent high-cost challenger prepare.
Take an incumbent distribution with a positive spread \(\DT\) and \(g_H>g_L\), and write \(\tau\) for its preparation
threshold.
\begin{enumerate}[label=(\alph*),leftmargin=2em]
\item A quiet market and an informative market coexist if and only if
\[
\pi<\tau\le M(\pi),\qquad \frac{k}{1-\tau}\le\DT\le\frac{k}{\rho\max\{\pi,1-\pi\}}.
\]
The informative market is then unique: class \(L\) abstains, only a purchase recruits a high-cost challenger, and
entry is
\[
\rho+(1-\rho)\min\Bigl\{\frac{(1-\pi)\,b\,\DT}{k},\;b+\pi\lambda\Bigr\},
\]
which rises strictly with the spread while class \(H\) mixes between a purchase and abstention and is constant once
class \(H\) buys for sure.
\item Along a path of such distributions, the strengths at which a quiet market coexists with an informative market
in which class \(H\) buys for sure form an interval, and every other strength with coexistence lies below it. Every
gap in coexistence lies at spreads too narrow for a sure purchase to pay.
\item Along a path on which the product of the spread and \(g_H-c_H\) does not fall with strength, the strengths
with coexistence form an interval, at every prior. Uniform incumbents on \([0,r]\), with \(r\) in a closed interval
above the reserve, satisfy this condition when every low-class value lies between the reserve and the lowest
strength, every high-class value is at least the highest strength, and class \(H\)'s mean value exceeds the highest
strength by at least \(c_H\).
\end{enumerate}
\end{paperresult}

In words: the threshold lies above the prior and within reach of a purchase; a purchase that just reaches it and
recruits covers \(k\); no order against a silent price does.

% =====================================================================
\section{The proof of part (a)}\label{sec:coexistence-proof-a}

If \(\pi\ge\tau\), a quiet market needs \(\DT\max\{\pi,1-\pi\}\le k\), and then every informed order earns less than
\(k\). So let \(\pi<\tau\). Appendix~A.4 splits every equilibrium into three pieces. By Lemma~A.3(b), a rate \(0\)
needs \(K\le k\), the full rate needs \(K\ge k\), and a mixed rate needs \(K=k\); each return falls strictly in its
rate.
\begin{itemize}
\item \emph{Lemma~A.7 (the high class).} The pair \((s_+,e(1))\) is one of five: no purchase (N); a mixed purchase
  below \(\tau\) at floor entry (Q); a mixed purchase above \(\tau\) that recruits, \(\DT[1-\mu_+(s_+)]=k\) (R); the
  threshold pair \(s_+=s^*\), \(e(1)=k/J\), which needs \(J=k\) when an indifferent challenger prepares (T); and a full
  purchase (F).
\item \emph{Lemma~A.8 (the low class).} A sale never recruits, so \(e(-1)=\rho\). When the investor may mix, the rate
  is \(0\) if \(\rho\DT\pi\le k\), \(\lambda\) if \(\rho\DT m\ge k\), and \(\hat s_-=b(\pi\rho\DT/k-1)/(1-\pi)\) in
  between. When it abstains, no rate qualifies if \(\rho\DT m\le k<\rho\DT\pi\).
\item \emph{Lemma~A.9 (product structure).} If \(s_+\ge s_-\), the zero order has a posterior at most \(\pi\), so
  \(e(0)=e(-1)=\rho\), and a pair of Lemma~A.7 and a rate of Lemma~A.8 form an equilibrium with the prices
  \(P(x)=t_0+e(x)w(\mu(x))\), \(w(\mu)=t_L-t_0+\DT\mu\).
\end{itemize}
A quiet market is (N) with \(s_-=0\), so it exists exactly when \(\DT\le\DT^{\rm N}\). Then \(s_-=0\). In an
informative market (Q) and (F) with \(M<\tau\) earn less than \(\rho\DT(1-\pi)\le k\), so the purchase recruits,
\(\tau\le M\), and \(k\le e(1)\DT[1-\mu_+(s_+)]\le J\). Conversely, if \(\DT\ge\DT^{\rm F}\), (F) works and is the only
pair. If \(\DT<\DT^{\rm F}\), the return falls continuously from \(J\ge k\) at \(s^*\) to \(\DT(1-M)<k\), so one rate
gives \(K_H=k\): (T) if \(J=k\), (R) otherwise. With \(\Pr(X=1)=b+\pi s_+\) and \(1-\mu_+(s_+)=k/\DT\),
\begin{equation}\label{eq:coexistence-entry}
\Ent=\rho+(1-\rho)\min\Bigl\{\frac{(1-\pi)\,b\,\DT}{k},\;b+\pi\lambda\Bigr\}.
\end{equation}
Indifference fixes the share of class \(L\) among purchases at \(k/\DT\). A wider spread needs more class-\(H\)
purchases next to the fixed liquidity pool, so entry rises.

% =====================================================================
\section{Parts (b) and (c): the path}\label{sec:coexistence-proof-b}

Along a path, \(\DT\) never falls and \(B_r(\mu)\) never rises. So \(\{\pi<\tau\}\) and \(\{\DT\ge\DT^{\rm F}\}\) are
up-sets, \(\{\tau\le M\}\) and \(\{\DT\le\DT^{\rm N}\}\) are down-sets, and their intersection is an interval. With a
sure purchase, \(J\ge\DT(1-M)\ge k\) holds by itself: part (b). Only \(J\ge k\) has no direction, and
\(J=\DT(g_H-c_H)/\gamma\) with \(\gamma\) never rising. If the numerator does not fall, \(\{J\ge k\}\) is an up-set:
part (c). For the uniform family, \(g_H-c_H\ge(r^2-p^2)/(2r)\) and \(r\DT'\ge\DT\) give
\(\bigl(\DT(g_H-c_H)\bigr)'\ge\frac{r^2-p^2}{2r^2}(r\DT'-\DT)\ge0\).

% =====================================================================
\section{Corollary 2: coexistence in the example}\label{sec:coexistence-example}

Take the example of Chapter~\ref{ch:example} at the prior \(1/2\), with a slot in the band of Corollary~1, so
\(1-M=m=1/(2+u)\). Each marked strength is where one order just earns \(k\): \(J(r_A)=k\);
\(\DT(r_F)(1-M)=k\); \(\rho\DT(r_N)/2=k\); \(\rho J(r_J)=k\); \(\rho\DT(r_S)m=k\); and \(\tau(r_C)=M\), the entry ceiling.
On the path, \(\tau\), \(J\) and \(\DT\tau\) rise strictly (Lemma~A.10), and Lemma~A.11 orders the strengths:
\[
r_0<r_A<r_F<r_N<r_J<r_S<r_1<r_C<r_2 .
\]
For example, \(r_F<r_N\) reads \(M<1-\rho/2\), true since \(M<\tau(r_2)<1-\rho\).

\begin{paperresult}{Corollary 2 (Coexistence in the example, the full statement of Appendix A.4)}
Take the example at the prior \(1/2\), a slot with \(0<\lambda<1\), \(0<\beta<1/2\) and \(\tau(r_1)<M(1/2)<\tau(r_2)\),
and the default tie rules. The equilibrium set at each strength of the path is
\begin{center}\upshape
\begin{tabular}{@{}ll@{}}
\toprule
Strength & Equilibria\\
\midrule
\(r_0\le r<r_A\) & \textbf{N}\\
\(r=r_A\) & \textbf{N}; \textbf{T} with entry one after a purchase\\
\(r_A<r<r_F\) & \textbf{N}; \textbf{R}\\
\(r_F\le r\le r_N\) & \textbf{N}; \textbf{F}\\
\(r_N<r<r_J\) & \textbf{F}; \textbf{Q}\\
\(r_J\le r\le r_2\) & \textbf{F}\\
\bottomrule
\end{tabular}
\end{center}
In each, the low class sells at the rate \(\hat s_-(r)\). Entry is \(\rho\) in a quiet market and in mixed orders at
the entry floor; \(\Ent_{\rm R}(r)=\rho+(1-\rho)\,b\DT(r)/(2k)\) in a mixed purchase that recruits, and
\(\Ent_{\rm R}(r_A)\) in the threshold pair at \(r_A\); and \(\Ent_{\rm F}=\rho+(1-\rho)(b+\lambda/2)\) in a full
purchase up to \(r_C\), with \(\Ent_{\rm F}=\rho\) above \(r_C\).
\begin{enumerate}[label=(\alph*),leftmargin=2em]
\item A quiet market and an informative market coexist exactly on \([r_A,r_N]\).
\item Along the informative branch, entry is \(\rho+(1-\rho)\min\{b\DT(r)/(2k),\,b+\lambda/2\}\) on \([r_A,r_C]\) and
\(\rho\) on \((r_C,r_2]\). It is continuous on \([r_A,r_C]\), rises strictly on \([r_A,r_F]\), is constant on
\([r_F,r_C]\), and returns to the floor \(\rho\) above the entry ceiling \(r_C\).
\item The equilibrium set is nonempty at every strength.
\end{enumerate}
\end{paperresult}

\paragraph{Proof.} At \(\pi=1/2\), \(1-\mu_+(s)=\mu_-(s)=b/(2b+s)\). So (N) holds for \(r\le r_N\); (R), with
\(s_R=b[\DT/k-2]\), on \((r_A,r_F)\); (F) for \(r\ge r_F\); (Q), with \(s_Q=b[\rho\DT/k-2]=\hat s_-\), on
\((r_N,r_J)\); and (T) at \(r_A\). Always \(s_+\ge s_-\), so Lemma~A.9 gives the table. Entry has no jump at \(r_F\),
since \(b+s_R/2=b\DT/(2k)=b+\lambda/2\) there.

\begin{figure}[htbp]
\centering
\includegraphics[width=0.6\linewidth]{coexistence}
\caption{Entry and the coexistence region along the example's path, at the prior \(1/2\), under the default tie
rules. Panels (a) and (c) plot entry against the strength \(r\) in each equilibrium of Corollary~2; panel (c)
enlarges the strengths around the shaded coexistence interval \([r_A,r_N]\). Panel (b) draws the region of
Proposition~3(a) in the plane of \(\DT\) and \(\tau\), with the curve \((\DT(r),\tau(r))\) of the path. The labels N,
R, F, Q and T name the quiet market, the mixed purchase that recruits, the full purchase, mixed orders at the entry
floor and the threshold pair. Chapter~\ref{app:numbers} lists the values of the marked strengths.}
\label{fig:coexistence-path}
\end{figure}

In Figure~\ref{fig:coexistence-path}, the informative branch rises along R from \(r_A\), stays on the plateau of F up
to \(r_C\), and drops to \(\rho\); the quiet market ends at \(r_N\).

\begin{watchout}
Proposition~2 holds under all four tie rules; Proposition~3 does not. If an indifferent investor abstains, R, T and Q
drop out, coexistence shrinks to \((r_F,r_N]\), and no equilibrium exists on \((r_N,r_S]\).
\end{watchout}

\section{Coexistence at other priors}\label{sec:coexistence-priors}

Corollary~A.2 of the paper takes any prior, any slot and the default tie rules. On the path, coexistence holds on
\[
\mathcal C(\pi)=\bigl\{r\in[r_0,r_2]:\ r\ge r_A,\ \ \pi<\tau(r)\le M(\pi),\ \ \rho\DT(r)\max\{\pi,1-\pi\}\le k\bigr\},
\]
because \(J\) rises, so \(J\ge k\) is \(r\ge r_A\), free of the prior and the slot. Let \(\tau_A=\tau(r_A)\), let
\(M(\pi_A^{\rm lo})=\tau_A\), and let \(\pi_{\max}=\tau(r_*)\), where \(\rho\DT(r_*)\tau(r_*)=k\). Then
\(\mathcal C(\pi)\) is empty for \(\pi<\pi_A^{\rm lo}\) (the top posterior cannot reach \(\tau_A\)); it starts at
\(r_A\) for \(\pi_A^{\rm lo}<\pi<\tau_A\); it starts just after the strength where \(\tau=\pi\) for
\(\tau_A\le\pi<\pi_{\max}\); and it is empty for \(\pi\ge\pi_{\max}\), where the quiet market ends first. Where it is not
empty, it is an interval that ends at the first strength where \(\tau\) passes \(M(\pi)\) or an order against a silent
price starts to pay.

\begin{chaptersummary}
A quiet and an informative market coexist exactly when \(\pi<\tau\le M(\pi)\) and
\(k/(1-\tau)\le\DT\le k/(\rho\max\{\pi,1-\pi\})\). Then class \(L\) abstains, only a purchase recruits, and entry
\(\rho+(1-\rho)\min\{(1-\pi)b\DT/k,\,b+\pi\lambda\}\) rises with the spread while class \(H\) mixes. Along a path the
sure-purchase strengths form one interval, and the balance condition makes the whole set one interval. In the example
at the prior \(1/2\), coexistence runs from \(r_A\) to \(r_N\), with \(r_0<r_A<r_F<r_N<r_J<r_S<r_1<r_C<r_2\).
\end{chaptersummary}
